\chapter{News and Events}\label{ch:s1:news-and-events}

A headline reaches the wire at 10:32:07. Within a few milliseconds, programs that parse the feed have traded on it, and most of the price's move is over
before a person has finished the first sentence. What is left is the part the market takes days to believe: Tetlock, Saar-Tsechansky and Macskassy found
that prices briefly underreact to the negative words in firm-specific news, and Hirshleifer, Lim and Teoh found that the reaction to an earnings surprise
is weaker, and the drift after it stronger, on days when many other firms announce. On this chapter's synthetic news stream a trader who acts within a
millisecond still has 96.7\% of the move ahead; at a tenth of a second, 70.1\%; after a second, only the 28.7\% that the machines leave for the following
days. A daily trader who reads each item's tone with an imperfect text model earns a Sharpe ratio of 3.2 on that remainder, more on busy news days than on
quiet ones. The build is \texttt{firm.newsevent}.

\section{Machine-readable news}

\begin{definition}[Machine-readable news]\label{def:s1:news-and-events:mrn}
\emph{Machine-readable news}\index{machine-readable news} is a news feed delivered in a structured form that programs can act on without a person reading
it: each item carries a timestamp, the securities it concerns, a category, and often scores for relevance, novelty and tone computed by the vendor or by
the user's own text model.
\end{definition}

A news strategy turns items into positions in three steps. It tags each item with the securities it concerns and discards those that are not about them.
It scores the item's direction and strength, from a dictionary of positive and negative words at the simplest to a language model trained on past
reactions. And it decides when to act: within the first second, when the move is still to come but only machines compete for it, or later, on the part
the market takes longer to absorb. The published record is mostly about the second kind. Tetlock measured the pessimism of a daily \emph{Wall Street
Journal} column and found that high pessimism predicts downward pressure on market prices followed by a reversion to fundamentals: tone as noise-trader
sentiment. Tetlock, Saar-Tsechansky and Macskassy took the measure to firm-specific stories and found that the fraction of negative words forecasts low
earnings and that prices briefly underreact to it, most for stories about fundamentals: tone as information that the market takes time to absorb.

The chapter's stream (\texttt{firm.newsevent}) has eight years of items on 1,000 stocks, 40.8 a day on average (82,351 in all), with a news intensity that
varies from day to day. Each item has a tone, the total move it will cause, with a standard deviation of 2\%; an attention level, which falls as more
items arrive the same day ($1/(1+0.03(k-1))$ for $k$ items, a median of 0.42); and a time of arrival, uniform over the session. A share of the move, from
one half with no attention to all of it with full attention, happens within a fraction of a second, with a time constant of 0.2 seconds; the rest drifts
in evenly over the next five days.

\section{Scheduled events}

Earnings announcements, dividend declarations, index changes and economic releases arrive at known times, and chapter~7 traded the drift after the first
of them with an \omterm{def:s1:earnings:event}{event study}. For news trading the schedule changes two things. The reaction can be prepared: the consensus is known, the surprise is a
single number, and the order can be staged before the release, which puts the first second in the hands of those with the fastest feeds. And the
distraction is predictable: Hirshleifer, Lim and Teoh measured the investor's information load by the number of same-day announcements and found that the
immediate price and volume reaction to a surprise is much weaker, and the post-announcement drift much stronger, when many other firms announce the same
day, strong enough to give a trading strategy substantial alphas. The synthetic stream plants their mechanism in its attention function.

\begin{definition}[News reaction window]\label{def:s1:news-and-events:window}
A \emph{news reaction window}\index{news reaction window} is the interval after a news item over which a strategy measures or trades its price move, from
the item's timestamp plus the strategy's latency to the end of the holding period; the share of the item's total move that falls inside it is what the
strategy can capture.
\end{definition}

\section{Unscheduled events and halts}

Unscheduled news (mergers, profit warnings, regulatory actions, accidents) arrives at any time, and the exchanges have a tool for the largest items: they
stop trading until the news is out.

\begin{definition}[Trading halt]\label{def:s1:news-and-events:halt}
A \emph{trading halt}\index{trading halt} is a stop in trading of a security ordered by its listing exchange or regulator, pending news, after news has
been released, or after a price move too large for the market to absorb continuously; trading resumes, usually through a reopening auction, when the
halt is lifted.
\end{definition}

\begin{dated}{2026-09}{Halt codes on Nasdaq}\label{dat:s1:news-and-events:halts}
Nasdaq publishes a code with each halt. Among them: T1, ``Halt -- News Pending''; T2, ``Halt -- News Released''; T3, news and resumption times; T5, a
single-stock trading pause; LUDP, a volatility trading pause under limit up--limit down; and MWC1 to MWC3, the three levels of the market-wide circuit
breaker. Book~1, chapter~31 describes the limit up--limit down bands and the circuit breakers.
\end{dated}

A halt changes who captures the move. While a stock is halted no one can trade it continuously, so the move happens at the reopening auction, where
everyone's orders meet at one price: speed within the first second is worth nothing, and what matters is how orders are placed in the auction and what
happens after it. The synthetic stream halts every item that moves its stock by more than 5\% (1.18\% of items) for five minutes, and puts the immediate
part of the move at the reopening. A fast trader catches nothing of a halted item's immediate move; the drift afterwards is the same as for any other item.
Whether reopening prices overshoot or undershoot is an empirical question the chapter's model does not settle; the post-halt strategy file sets out what
to test.

\section{Reaction speed and who captures the move}

\Cref{lst:s1:news-and-events:remaining} is the whole reaction model: the share of an item's move still ahead of a trader who acts at a given latency.
Weighting the items by the size of their moves gives the share of the stream's total move that is still available.

\begin{center}
\small
\begin{tabular}{@{}lcccccc@{}}
\toprule
latency & 1 ms & 10 ms & 100 ms & 1 s & 1 hour & next close\\
\midrule
share of the move still ahead & 96.7\% & 93.7\% & 70.1\% & 29.2\% & 28.7\% & 28.7\%\\
P\&L per item after 10 bp (bp) & 144 & 139 & 102 & 36 & 36 & 36\\
\bottomrule
\end{tabular}
\end{center}

The curve (\cref{fig:s1:news-and-events:latency}) has two regimes. Below a second the share falls fast, because the machines take the immediate part
within a few time constants: a trader a tenth of a second late has lost nearly a third of the move to faster ones. Above a second it is flat: the machines are
done, and what remains, 28.7\%, is the part the market's attention leaves for the next days. The P\&L figures assume the trader knows the direction of
every item, as the fastest traders effectively do for the simplest items; they are the value of the window, not of a strategy.

\begin{omfigure}
\begin{tikzpicture}
\begin{axis}[width=10cm,height=5.4cm,xmode=log,xlabel={\small latency after the item (seconds)},ylabel={\small share of the move still ahead (\%)},
  tick label style={font=\footnotesize},xmin=0.001,xmax=100,ymin=0,ymax=100,grid=major,grid style={gray!25},legend pos=north east,
  legend style={font=\scriptsize},table/col sep=comma]
\addplot[omThm,very thick] table[x=latency,y=share_pct]{figdata/strategies-1/17-news-and-events/latency.csv};
\addplot[black!40,dashed] coordinates {(0.001,28.7) (100,28.7)};
\legend{synthetic news stream,left at the next close}
\end{axis}
\end{tikzpicture}

\omcaption{The share of the synthetic news stream's total move still ahead of a trader who acts at a given latency after each item, weighted by the size
of each item's move. Machines take the immediate part within a fraction of a second; the dashed line is the drift left at the next close. Data:
\texttt{s1\_news.capture}.}\label{fig:s1:news-and-events:latency}
\end{omfigure}

A daily trader cannot compete in the first second, but has the rest. The chapter's daily trader reads each item's tone with a text model whose error is
1.5 times the tone's spread, so it gets the direction right 68.9\% of the time; it buys or sells at the next close in the direction of its reading and
holds five days, with 2\% of daily noise on each stock and ten basis points a round trip (\cref{lst:s1:news-and-events:daily}).

\begin{center}
\small
\begin{tabular}{@{}lccc@{}}
\toprule
items & all & quiet days (attention above median) & busy days\\
\midrule
items traded & 82,351 & 39,735 & 42,616\\
drift captured per item (bp) & 25.4 & 20.8 & 29.5\\
return a year after costs & 7.4\% & 6.3\% & 8.3\%\\
Sharpe ratio after costs & 3.24 & 1.80 & 2.74\\
\bottomrule
\end{tabular}
\end{center}

Busy days leave more drift per item, as in Hirshleifer, Lim and Teoh: the market's attention is spread over more news, so less of each move happens at
once. The split books each hold about half the items, so their Sharpe ratios are lower than the whole book's, which diversifies over twice as many. The
Sharpe ratio of 3.2 is high because the model plants the underreaction cleanly and the reading's errors are independent of everything else; the published
effects are measured in basis points per story, and the costs of trading many small positions quickly absorb them.

\section{Strategy files}

\begin{strategyfile}{Headline sentiment}\label{strat:s1:news-and-events:sentiment}
\sfield{Who pays you, and why} Investors who absorb the information in news text slowly, and noise traders who push prices on tone alone.
\sfield{Instruments and venues} Liquid stocks with frequent firm-specific news; index futures for market-wide tone.
\sfield{Signal} The tone of firm-specific stories from a dictionary or a trained text model, filtered for relevance and novelty.
\sfield{Sizing and execution} Positions after the first seconds, in the direction of the tone, held days; many small positions.
\sfield{Costs} News feeds and text processing; turnover on every story.
\sfield{How it dies} Better text models everywhere; the drift absorbed faster.
\sfield{Horizon, capacity, infrastructure} Days; capacity limited by the number of stories and their liquidity; a timestamped news archive.
\sfield{Backtest honestly} Stories as they arrived, with the feed's own timestamps, not the article's publication time; no revised tone scores.
\sfield{Sources} Tetlock (2007); Tetlock, Saar-Tsechansky and Macskassy (2008): prices briefly underreact to negative words.
\end{strategyfile}

\begin{strategyfile}{Scheduled-event drift}\label{strat:s1:news-and-events:scheduled}
\sfield{Who pays you, and why} Investors who underreact to scheduled news, by more when many announcements compete for their attention.
\sfield{Instruments and venues} Stocks with announcements on the calendar.
\sfield{Signal} The surprise against consensus, read from the release as it arrives.
\sfield{Sizing and execution} Positions at the close of the announcement day, in the direction of the surprise, held days to weeks.
\sfield{Costs} Moderate; trading near crowded closes.
\sfield{How it dies} Faster reactions by machines leave less drift.
\sfield{Horizon, capacity, infrastructure} Days to weeks; an event calendar and a consensus database.
\sfield{Backtest honestly} Release times to the second; consensus as it stood before the release.
\sfield{Sources} Chapter~7's \omterm{def:s1:earnings:event}{event study}; Hirshleifer, Lim and Teoh (2009).
\end{strategyfile}

\begin{strategyfile}{Post-halt reopening}\label{strat:s1:news-and-events:halt}
\sfield{Who pays you, and why} Traders who must trade at the reopening whatever the price, and a reopening auction whose price may overshoot.
\sfield{Instruments and venues} Stocks halted for news or volatility; the reopening auction.
\sfield{Signal} The news released during the halt, the auction's indicative price and imbalance, and the move since the halt.
\sfield{Sizing and execution} Orders in the reopening auction or just after it, against an overshoot or with a drift; small, because liquidity is thin.
\sfield{Costs} Wide spreads after the reopening; the risk of another halt.
\sfield{How it dies} Competition for reopening liquidity.
\sfield{Horizon, capacity, infrastructure} Minutes to days; the halt feed and the auction's indicative prices.
\sfield{Backtest honestly} Halts and resumption times from the exchange's records; auction prices, not the first trades after.
\sfield{Sources} No performance figure verified; the halt codes of the dated box.
\end{strategyfile}

\begin{strategyfile}{Underreaction to low-attention news}\label{strat:s1:news-and-events:attention}
\sfield{Who pays you, and why} Distracted investors, who react less to news on days when more other news competes for their attention.
\sfield{Instruments and venues} Stocks with news on busy days.
\sfield{Signal} The item's direction, weighted by the day's news load (the number of same-day announcements or stories).
\sfield{Sizing and execution} Larger positions on busy days; held days.
\sfield{Costs} As for headline sentiment.
\sfield{How it dies} Machines that do not get distracted.
\sfield{Horizon, capacity, infrastructure} Days; a count of the day's news.
\sfield{Backtest honestly} The news load as it was known at the time of the trade.
\sfield{Sources} Hirshleifer, Lim and Teoh (2009): stronger drift and substantial alphas; this chapter's split by attention.
\end{strategyfile}

\section{Tutorial: before a person reads it}\label{tut:s1:news-and-events:tut}

\begin{tutorial}
\textbf{Goal.} Generate a news stream with planted tone and attention, measure the share of its move still ahead at each latency, and trade what is left
at the next close. \textbf{End state:} the two tables and \cref{fig:s1:news-and-events:latency}.
\begin{tutsteps}
\item \textbf{The reaction model}: the immediate share with attention, and what is left at a given latency; halted items have nothing to catch.
\omcode{code/firm/newsevent/firm_newsevent.py}{55}{68}{The immediate share and the move still ahead at a latency.}\label{lst:s1:news-and-events:remaining}
\item \textbf{The daily trader}: a noisy reading of each item's tone, entry at the next close, five days' holding.
\omcode{code/strategies-1/17-news-and-events/python/s1_news.py}{48}{68}{The daily trader's book by attention.}\label{lst:s1:news-and-events:daily}
\item \textbf{Run} \texttt{stream()}, \texttt{capture(latency)} for each latency and \texttt{daily(split)} for the three splits, and \texttt{fig\_news.py}.
\end{tutsteps}
\textbf{What to change next.} Let the text model's error vary with the item's category; make the reopening auction overshoot and trade against it; give
the daily trader a signal weighted by the day's news load.
\end{tutorial}

\section{Build: a news stream}\label{bld:s1:news-and-events:newsevent}

\begin{build}
\bfield{Purpose} A synthetic news stream with tone, attention and halts, and the share of each item's move still ahead at a given latency.
\bfield{Interface} \texttt{NewsConfig(\dots)}, \texttt{simulate\_news(days, stocks, cfg, rng)}, \texttt{immediate\_share(attention, cfg)},
\texttt{remaining(items, latency, cfg)} (latency \texttt{None} for the next close).
\bfield{Rules} Attention depends only on the day's number of items; halted items reopen at the new price; the stream is deterministic for a given seed.
\bfield{Acceptance tests} \texttt{code/firm/newsevent/tests/}: the halt threshold and the attention function; the remaining share at zero, small and large
latencies.
\bfield{Stretch} Relevance and novelty scores; items that concern several stocks; reopening auctions with an imbalance.
\end{build}

\begin{omsources}
\item P.~C.~Tetlock, ``Giving content to investor sentiment: the role of media in the stock market'', \emph{Journal of Finance} 62(3), 2007.
\item P.~C.~Tetlock, M.~Saar-Tsechansky and S.~Macskassy, ``More than words: quantifying language to measure firms' fundamentals'', \emph{Journal of
Finance} 63(3), 2008.
\item D.~Hirshleifer, S.~S.~Lim and S.~H.~Teoh, ``Driven to distraction: extraneous events and underreaction to earnings news'', \emph{Journal of
Finance} 64(5), 2009.
\end{omsources}

\section{Exercises}

\begin{exercise}[$\star$]\label{exo:s1:news-and-events:1}
The immediate part of a move decays with a time constant of 0.2 seconds. What fraction of it is still ahead of a trader at 100 milliseconds, and at one
second?
\end{exercise}

\begin{exercise}[$\star$]\label{exo:s1:news-and-events:2}
On a day with 40 items, what is each item's attention, and what share of its move happens at once? What is the share at the median attention of 0.42?
\end{exercise}

\begin{exercise}[$\star$]\label{exo:s1:news-and-events:3}
Why does a fast trader catch nothing of a halted item's immediate move?
\end{exercise}

\begin{exercise}[$\star\star$]\label{exo:s1:news-and-events:4}
A text model reads a normally distributed tone with an independent normal error 1.5 times the tone's standard deviation. How often does it get the
direction right? (For a bivariate normal pair with correlation $\rho$, the signs agree with probability $\tfrac12 + \arcsin(\rho)/\pi$.)
\end{exercise}

\begin{exercise}[$\star\star$]\label{exo:s1:news-and-events:5}
Why is the latency curve flat between one second and the next close?
\end{exercise}

\begin{exercise}[$\star\star$]\label{exo:s1:news-and-events:6}
Why do the quiet-day and busy-day books both have lower Sharpe ratios than the whole book?
\end{exercise}

\begin{exercise}[$\star\star\star$]\label{exo:s1:news-and-events:7}
\emph{Coding.} Set \texttt{READ} to 0 in \texttt{s1\_news.py}, a perfect reading, and rerun \texttt{daily()}. What happens to the Sharpe ratio, and what does
the difference measure?
\end{exercise}

\begin{exercise}[$\star\star\star$]\label{exo:s1:news-and-events:8}
\emph{Find the flaw.} ``Our backtest trades each story at the minute of its publication time on the website and earns 40 basis points per story.''
\end{exercise}

\section{Problem: Before a Person Reads It}

\begin{problem}[{Weekend problem --- the life of a headline}]\label{pb:s1:news-and-events:1}
The chapter's synthetic news stream and the public record.

\medskip\noindent\textbf{Part I --- News as data.}
\begin{enumerate}
\item Define \omterm{def:s1:news-and-events:mrn}{machine-readable news} and a \omterm{def:s1:news-and-events:window}{news reaction window}.
\item List the three steps from an item to a position.
\item What did Tetlock (2007) find about media pessimism?
\item What did Tetlock, Saar-Tsechansky and Macskassy find about negative words?
\end{enumerate}

\medskip\noindent\textbf{Part II --- Events and halts.}
\begin{enumerate}[resume]
\item What changes when an event is scheduled?
\item What did Hirshleifer, Lim and Teoh find about distraction?
\item Define a \omterm{def:s1:news-and-events:halt}{trading halt} and name three Nasdaq halt codes.
\item How does a halt change who captures the move?
\end{enumerate}

\medskip\noindent\textbf{Part III --- The stream.}
\begin{enumerate}[resume]
\item Describe the synthetic stream: items, tone, attention, reaction and halts.
\item How many items are there, how many a day, and what share is halted?
\item Give the share of the move still ahead at 1 millisecond, 100 milliseconds, one second and the next close.
\item Explain the two regimes of the latency curve.
\end{enumerate}

\medskip\noindent\textbf{Part IV --- The verdict.}
\begin{enumerate}[resume]
\item State the \emph{named result}: the share of the news move captured at each latency, and the drift left for a daily trader.
\item Why do busy days leave more drift?
\item What does the text model's error cost the daily trader?
\item Why is the daily trader's Sharpe ratio of 3.2 an upper bound for real news?
\item How would you backtest a news strategy honestly?
\item Which strategy file needs the fastest infrastructure?
\item How does this chapter relate to chapter~7?
\item In one sentence: what does speed buy in news trading?
\end{enumerate}
\end{problem}

\section{Interview questions}

\begin{interviewq}[$\star$ \iqroles{researcher}]\label{iq:s1:news-and-events:1}
How would you turn a news feed into a trading signal?
\end{interviewq}

\begin{interviewq}[$\star\star$ \iqroles{researcher}]\label{iq:s1:news-and-events:2}
Your news backtest shows strong returns. What timestamps would you check first?
\end{interviewq}

\begin{interviewq}[$\star\star$ \iqroles{trader}]\label{iq:s1:news-and-events:3}
A stock you hold is halted with news pending. What do you do?
\end{interviewq}

\begin{interviewq}[$\star\star$ \iqroles{developer}]\label{iq:s1:news-and-events:4}
Design the path from a news vendor's feed to an order, for a strategy that must act within ten milliseconds.
\end{interviewq}

\begin{interviewq}[$\star\star$ \iqroles{researcher, trader}]\label{iq:s1:news-and-events:5}
Why might a news signal work better on days with many earnings announcements?
\end{interviewq}

\begin{interviewq}[$\star\star\star$ \iqroles{researcher}]\label{iq:s1:news-and-events:6}
An item's immediate move decays as $e^{-t/\tau}$ and a share $a$ of the total move is immediate. Derive the share of the move still ahead of a trader at
latency $L$, and the latency at which half of the immediate part is gone.
\end{interviewq}
